\subsection{A proposed causal experiment}
\label{sec:causal-experiment}

A fixed-workload experiment would compare teams given a specified set of AI tools and working practices with teams using the same ordinary development and review process without generative AI. Before assigning teams to groups, fix each team's task list, requirement versions, acceptance criteria and reference scores. The comparison then covers work specified before AI access is assigned. Pair teams by repository domain, previous workload and experience. Within each pair, use an independent fair coin to choose which team receives AI access \cite{imai2009}. Teams remain in their assigned groups throughout the study. Compare them according to this assignment regardless of actual tool use. Use task diaries or tool-usage logs to describe which tools they actually used.

The team boundaries also need to reflect how work is shared. Shared reviewers, developers, computing capacity or code can make one team's labour or requirement fulfilment depend on another team's AI access \cite{hudgens2008}. If these links require allocating several teams together, revise the estimator weights and planning before randomisation so that each team still receives equal weight in the comparison.

Each team's initial reference scores must have a positive total. Keep that total unchanged and use it as the denominator for both labour through task age $\tau+\Delta$ and the sum of output scores. Include labour on failed and abandoned tasks, and give an output score of zero to tasks that miss the acceptance deadline. The resulting team outcomes are labour per reference unit and output score per reference unit. The primary effects are the mean differences between groups in these outcomes, with equal weight for each team. Estimate each effect by subtracting the comparator team's outcome from the AI team's within each pair, then averaging the pair differences.

A benefit claim requires both a worthwhile labour reduction and an output loss smaller than a maximum set before inspecting the teams' labour and output scores. Evidence that fails to detect a quality reduction can still allow a loss exceeding that maximum \cite{piaggio2012}.

The provisional criteria require a saving exceeding one person-hour per reference unit and a loss of less than five percentage points in output score per reference unit. For uncertainty and measurement bias, the protocol uses approximate one-sided 97.5\% Student-t bounds from independent pair differences, expanded by 0.2 person-hours per reference unit for labour and 0.01 for output score per reference unit. Both expanded bounds must meet their respective criteria.\footnote{The Student $t$ distribution accounts for estimating the standard deviation from the sample. With $p$ pairs, an upper bound for the mean difference is the sample mean plus $t_{p-1,0.975}$ times its estimated standard error. Here $p-1$ is the degrees of freedom and $t_{p-1,0.975}$ is the distribution's 97.5th percentile. Exact Student $t$ coverage assumes independent, identically distributed normal pair differences. The protocol uses this reference distribution as an approximation. The added allowances widen the bounds for measurement bias and need separate empirical justification.} The measurement pilot must test whether time measurements and requirement assessments justify these bias allowances. The practical thresholds must be agreed before the main study. This plan permits a benefit claim only when labour per reference unit and output score per reference unit are available for every team.

The accompanying protocol specifies a separate six-pair measurement pilot. The acceptance deadline and the delay from acceptance to assessment are provisionally set to $\tau=\Delta=336$ hours (14 calendar days each). Twenty thousand simulated experiments per design and scenario compare 48 and 96 teams using hypothetical distributions of labour and output scores, measurement bias and missing team measurements, including cases at each benefit boundary \cite{morris2019}. In the required hypothetical scenarios with a worthwhile AI benefit, neither team count achieves the required 90\% probability of meeting both criteria.

Before choosing the main study's sample size, the pilot must establish defensible ranges for measurement error and variation between teams. It must also establish how often both the labour and output measures can be obtained for a team. The pilot's estimated mean effect of AI access must not determine that choice \cite{albers2018}. The experiment is specified but has not been run. The simulations concern hypothetical labour and output scores rather than measured workplace effects.
